
\subsection{First Launch}
When you first launch the app, you will be prompted to grant permissions to the app.

\begin{figure}[H]
    \centering
    % Left Minipage
    \begin{minipage}[t]{0.22\textwidth}
      \centering
      \includegraphics[width=\linewidth]{Walthrough/Simulator Screenshot - iPhone 17 Pro - 2026-03-19 at 15.07.11.png}
      \caption{\makebox[\linewidth][c]{App installed on phone}}
      \label{fig:Installed}
    \end{minipage}
    \quad
    % Right Minipage
    \begin{minipage}[t]{0.22\textwidth}
      \centering
      \includegraphics[width=\linewidth]{Walthrough/Simulator Screenshot - iPhone 17 Pro - 2026-03-19 at 15.07.35.png}
      \caption{Welcome screen}
      \label{fig:FirstLaunchWelcome}
    \end{minipage}
  \end{figure}
Select \textbf{Allow while using app} for location and \textbf{Allow} for activity detection for motion.

\begin{figure}[H]
    \centering
    \begin{minipage}[t]{0.22\textwidth}
      \centering
      \includegraphics[width=\linewidth]{Walthrough/Simulator Screenshot - iPhone 17 Pro - 2026-03-19 at 15.10.37.png}
      \caption{Location permission request}
      \label{fig:FirstLaunchPermissionsLocation}
    \end{minipage}
    \quad
    \begin{minipage}[t]{0.22\textwidth}
      \centering
      \includegraphics[width=\linewidth]{Walthrough/Simulator Screenshot - iPhone 17 Pro - 2026-03-19 at 15.10.45.png}
      \caption{Allow location access}
      \label{fig:FirstLaunchPermissionsLocationAccept}
    \end{minipage}
    \quad
    \begin{minipage}[t]{0.22\textwidth}
      \centering
      \includegraphics[width=\linewidth]{Walthrough/Simulator Screenshot - iPhone 17 Pro - 2026-03-19 at 15.10.51.png}
      \caption{Motion permission request}
      \label{fig:FirstLaunchPermissionsMotion}
    \end{minipage}
    \quad
    \begin{minipage}[t]{0.22\textwidth}
      \centering
      \includegraphics[width=\linewidth]{Walthrough/Simulator Screenshot - iPhone 17 Pro - 2026-03-19 at 15.11.36.png}
      \caption{Allow motion access}
      \label{fig:FirstLaunchPermissionsMotionAccept}
    \end{minipage}
\end{figure}

\subsubsection{Setting your device ID}
You will need to enter your unique device ID which I will provide when we meet. Tap the arrow to proceed.
\begin{figure}[H]
    \centering
    \begin{minipage}[t]{0.22\textwidth}
      \centering
      \includegraphics[width=\linewidth]{Walthrough/Simulator Screenshot - iPhone 17 Pro - 2026-03-19 at 15.11.06.png}
      \caption{Enter device ID}
      \label{fig:DeviceIDblank}
    \end{minipage}
    \quad
    \begin{minipage}[t]{0.22\textwidth}
      \centering
      \includegraphics[width=\linewidth]{Walthrough/Simulator Screenshot - iPhone 17 Pro - 2026-03-19 at 15.11.25.png}
      \caption{Device ID entered}
      \label{fig:DeviceIDfilled}
    \end{minipage}
\end{figure}
\subsection{Tracking tab}
This is the main tab of the app. It will show you: 
\begin{itemize}
    \item{Your current location}
    \item {An option to temporarily pause the tracking, if you want to remain private for a short period of time.}
    \item {Your device ID}
    \item {The status of your home location (set or not set)}
\end{itemize}

\subsubsection{Start sharing your location}
You will need to start sharing your location. Tap the toggle to proceed.

\begin{figure}[H]
    \centering
    \begin{minipage}[t]{0.22\textwidth}
      \centering
      \includegraphics[width=\linewidth]{Walthrough/Simulator Screenshot - iPhone 17 Pro - 2026-03-19 at 15.13.14.png}
      \caption{Turn on location sharing}
      \label{fig:StartSharingLocation}
    \end{minipage}
    \quad
    \begin{minipage}[t]{0.22\textwidth}
      \centering
      \includegraphics[width=\linewidth]{Walthrough/Simulator Screenshot - iPhone 17 Pro - 2026-03-19 at 15.13.21.png}
      \caption{Set home location prompt}
      \label{fig:SettingHomeLocation1}
    \end{minipage}
    \end{figure}


\subsubsection{Setting your home location}
You will need to set your home location. You'll only need to do this once. Ideally, do this the first time you are at home after installing the app. Tap the arrow to proceed.

\begin{figure}[H]
    \centering
    \begin{minipage}[t]{0.22\textwidth}
    \centering
    \includegraphics[width=\linewidth]{Walthrough/Simulator Screenshot - iPhone 17 Pro - 2026-03-19 at 15.13.21.png}
    \caption{\makebox[\linewidth][c]{Select home location}}
    \label{fig:SettingHomeLocation2}
  \end{minipage}
  \quad
  \begin{minipage}[t]{0.22\textwidth}
    \centering
    \includegraphics[width=\linewidth]{Walthrough/Simulator Screenshot - iPhone 17 Pro - 2026-03-19 at 15.13.34.png}
    \caption{Confirm home location}
    \label{fig:SettingHomeLocation3}
  \end{minipage}
  \end{figure}

\footnote{This location will never be stored or shared with anyone. It will be used as the origin of a relative co-ordinate system for your location data. For example, if you visit a cafe it may be 4600m north-west of your home. The exact location or name of the cafe will not be stored.}

\subsubsection{Pausing location sharing}
You can pause location sharing for 30 minutes by toggling off location sharing. This is to prevent the app from using excessive battery life or to maintain privacy should you wish to remain private for a short period of time. Tap resume now to start location sharing again when ready.

\begin{figure}[H]
    \centering
    \begin{minipage}{0.22\textwidth}
      \centering
      \includegraphics[width=\linewidth]{Walthrough/Simulator Screenshot - iPhone 17 Pro - 2026-03-19 at 15.15.59.png}
      \caption{Pause location sharing}
      \label{fig:PausingLocationSharing}
    \end{minipage}
  \end{figure}




\subsection{Diary tab}
Diaries are completed on days \textbf{2-11}, for the previous day.
\subsubsection*{Selecting a day}
Select yesterday to complete diary the previous day.

If you forget, select another date to complete the diary for. In this case it is a good idea to refresh the diary before completing it.
\begin{figure}[H]
    \centering
    \begin{minipage}{0.22\textwidth}
      \centering
      \includegraphics[width=\linewidth]{Walthrough/Simulator Screenshot - iPhone 17 Pro - 2026-03-19 at 15.21.18.png}
      \caption{Select diary day}
      \label{fig:SelectingADay}
    \end{minipage}
  \end{figure}



\subsubsection*{Filling in and submitting the diary}
To complete a diary, review the day's entries and journeys.

Some entries will match what you did that day. For example, `0900-1000: you were at home (residing there)`. In these cases, confirm the place and what you were doing by choosing the appropriate options from the dropdown menus. If you need to, you can add extra context in the free text box (for example, "I was at home, doing some work and watching a movie").

Other entries will not match what you did. For example, `1000-1005: it suggests you were at a waterpark relaxing`. In these cases, reject the suggested activity/place using the dropdown options. Before you can submit the diary, you must add extra context in the free text box (for example, "I was at the waterpark, but I was working").

In a similar way, you will be asked about your journeys that day, and to confirm the transport you used for each journey.

Once opened, you can edit each diary as many times as you like until you submit it.

\begin{figure}[H]
    \centering
    \begin{minipage}[t]{0.22\textwidth}
      \centering
      \includegraphics[width=\linewidth]{Walthrough/Simulator Screenshot - iPhone 17 Pro - 2026-03-09 at 07.52.10.png}
      \caption{Fill visit entry}
      \label{fig:FillingAVisit}
    \end{minipage}
    \quad
    \begin{minipage}[t]{0.22\textwidth}
      \centering
      \includegraphics[width=\linewidth]{Walthrough/Simulator Screenshot - iPhone 17 Pro - 2026-03-09 at 07.53.14.png}
      \caption{Fill journey entry}
      \label{fig:FillingAJourney}
    \end{minipage}
  \end{figure}

Once a diary is complete, you can press the submit button to send it to the database. You can only submit once. After you submit, your answers can't be changed, but the app will warn you before you confirm.

\Callout{Diary Advice}{%
The place category is broad (for example, `restaurant` for a place that serves food and drink). The place type tries to narrow it down (for example, `Italian restaurant`). The activity label is a simple mapping (for example, at a restaurant: \textbf{Eating/Drinking}).

If the suggestion is mostly right, confirm the entry and add extra context if you need to. If it isn't right, reject the relevant part and add extra context about where you actually were and what you were doing.}

\subsection{Troubleshooting}
If you are having trouble with the app, please contact me at \href{mailto:alex.m.skondras@kcl.ac.uk}{alexander.m.skondras@kcl.ac.uk}.